\begin{center}[18.]\end{center}
Fügt man zu \textit{den Massen im Inneren eines Körperraums} auf der Oberfläche die Massen mit entgegengesetztem Vorzeichen hinzu, die ihnen für den Außenraum gleichwertig sind, so erhält man einen Körper, der positive und negative Massen trägt. Zusammen üben diese auf keinen Punkt des Außenraums eine Anziehungskraft aus.

\SourceLeaf{629}
Man kann leicht beweisen

1.\quad dass dieser Körper auch infolge der Reaktion äußerer Massen im Gleichgewicht bleibt

2.\quad dass der Körper auch auf einen Rheophor keinerlei Kraft ausübt, wenn die betreffenden Massen magnetische Fluida sind.

Schwerer ist dagegen zu beweisen

3.\quad dass dieser Körper auch trotz der Reaktion des Rheophors im Gleichgewicht bleibt.

Der letzte Punkt beruht auf den folgenden Überlegungen.

Sei $\mathrm dm$ ein Element des Körpers mit den Koordinaten $x,y,z$. Das Zeichen $\int$ bezieht sich auf alle $\mathrm dm$. Sei ferner $\mathrm ds$ ein Element eines galvanischen Stroms mit den Endpunktkoordinaten $a,b,c;\ a+\mathrm da,\ b+\mathrm db,\ c+\mathrm dc$. Also sind $a,b,c$ Funktionen von $s$. Die Stromstärke ist $=1$. Das Zeichen $S$ bezeichnet die Summation in Bezug auf $\mathrm ds$.

Sind $X,Y,Z$ die Komponenten der gesamten auf $\mathrm dm$ wirkenden beschleunigenden Kraft, so lauten die Gleichgewichtsbedingungen bekanntlich
\[
\begin{gathered}
\int X\mathrm dm,\int Y\mathrm dm,\int Z\mathrm dm,\int(zY-yZ)\mathrm dm,\\
\int(xZ-zX)\mathrm dm,\int(yX-xY)\mathrm dm\text{ alle }=0
\end{gathered}
\]

Es sei $(x-a)^2+(y-b)^2+(z-c)^2=rr$ (die positiven $x,y,z$ sind nach vorn, rechts beziehungsweise oben gerichtet).
\[
\begin{aligned}
X&=S\frac{(y-b)\mathrm dc-(z-c)\mathrm db}{r^3},\\
Y&=S\frac{(z-c)\mathrm da-(x-a)\mathrm dc}{r^3},\\
Z&=S\frac{(x-a)\mathrm db-(y-b)\mathrm da}{r^3}
\end{aligned}
\]
\[
\text{also}\qquad\int X\mathrm dm=S(\eta\mathrm dc-\zeta\mathrm db)
\]
\[
\begin{aligned}
&\int(yX-xY)\mathrm dm\\
&=SV\mathrm dc+S\{(\xi c-\zeta a)\mathrm da+(\eta c-\zeta b)\mathrm db+(a\xi+b\eta+c\zeta)\mathrm dc\}\\
&\qquad-\int z\mathrm dmS\frac{\mathrm d\frac1r}{\mathrm ds}\mathrm ds
\end{aligned}
\]
[Dabei ist $\displaystyle\int\frac{\mathrm dm}{r}=V,\quad\int\frac{x-a}{r^3}\mathrm dm=\xi$ usw. gesetzt. Alle diese Größen werden null.]

\medskip
1836 Februar 18.

Die Ableitung lässt sich auf folgende Weise viel einfacher führen.

Sind $x,y,z;\ x+\mathrm dx,\ y+\mathrm dy,\ z+\mathrm dz$ die Koordinaten zweier beliebiger, einander unendlich naher Punkte, so muss die Variation von
\[
\mathrm dx^2+\mathrm dy^2+\mathrm dz^2
\]
\SourceLeaf{630}
unabhängig von den Werten von $\mathrm dx,\mathrm dy,\mathrm dz,$ $=0$ werden. Zur Abkürzung bezeichnen wir $\delta x,\delta y,\delta z$ mit $\alpha\xi,\alpha\eta,\alpha\zeta$, wobei $\alpha$ einen konstanten, unendlich kleinen Koeffizienten bezeichnet. Also gilt
\[
\begin{aligned}
\delta(\mathrm dx^2+\mathrm dy^2+\mathrm dz^2)&=2\alpha(\mathrm dx.\mathrm d\xi+\mathrm dy.\mathrm d\eta+\mathrm dz.\mathrm d\zeta)\\
&=2\alpha\Big\{\frac{\mathrm d\xi}{\mathrm dx}.\mathrm dx^2+\frac{\mathrm d\eta}{\mathrm dy}.\mathrm dy^2+\frac{\mathrm d\zeta}{\mathrm dz}.\mathrm dz^2\\
&\quad+\left(\frac{\mathrm d\xi}{\mathrm dy}+\frac{\mathrm d\eta}{\mathrm dx}\right)\mathrm dx\mathrm dy+\left(\frac{\mathrm d\xi}{\mathrm dz}+\frac{\mathrm d\zeta}{\mathrm dx}\right)\mathrm dx\mathrm dz\\
&\quad+\left(\frac{\mathrm d\eta}{\mathrm dz}+\frac{\mathrm dz}{\mathrm dy}\right)\mathrm dy\mathrm dz\Big\}
\end{aligned}
\]
Also muss offenbar Folgendes gelten.
\[
\begin{aligned}
1.\quad0&=\frac{\mathrm d\xi}{\mathrm dx}\\
2.\quad0&=\frac{\mathrm d\eta}{\mathrm dy}\\
3.\quad0&=\frac{\mathrm d\zeta}{\mathrm dz}\\
4.\quad\frac{\mathrm d\xi}{\mathrm dy}&=-\frac{\mathrm d\eta}{\mathrm dx}=r\\
5.\quad\frac{\mathrm d\zeta}{\mathrm dx}&=-\frac{\mathrm d\xi}{\mathrm dz}=q\\
6.\quad\frac{\mathrm d\eta}{\mathrm dz}&=-\frac{\mathrm d\zeta}{\mathrm dy}=p
\end{aligned}
\]
Aus (6) (4) (5) (6), aus (6) (2) und aus (6) (3) erhält man
\[
\begin{aligned}
\frac{\mathrm dp}{\mathrm dx}&=\frac{\mathrm dd\eta}{\mathrm dx\mathrm dz}=-\frac{\mathrm dd\xi}{\mathrm dz\mathrm dy}=\frac{\mathrm dd\zeta}{\mathrm dy\mathrm dx}=-\frac{\mathrm dp}{\mathrm dx}=0\\
\frac{\mathrm dp}{\mathrm dy}&=\frac{\mathrm dd\eta}{\mathrm dy\mathrm dz}=0\\
\frac{\mathrm dp}{\mathrm dz}&=-\frac{\mathrm dd\zeta}{\mathrm dz\mathrm dy}=0
\end{aligned}
\]
Also ist $p$ konstant.

Ebenso folgt, dass $q,r$ konstant sind,

und daher gilt
\[
\begin{aligned}
\xi&=a+ry-qz\\
\eta&=b+pz-rx\\
\zeta&=c+qx-py
\end{aligned}
\]

\begin{center}[19.]\par\medskip\textit{Beweis von} \textsc{Ampères} \textit{Fundamentalsatz.}\end{center}
Über eine begrenzte Fläche (I), auf der ein beliebiger Punkt die Koordinaten $x,y,z$ hat, sei positives magnetisches Fluidum gleichmäßig so verteilt, dass die Menge des magnetischen Fluidums pro Flächeneinheit $=k$ ist.

Ein Element eines galvanischen Stroms der Stärke $i$ erstrecke sich von $0,0,0$ bis $0,0,\zeta$.

\SourceLeaf{631}
Zur Abkürzung schreibt man
\[
r=\BareRoot{(xx+yy+zz)}
\]

Einem Element $\omega$ der $x,y$-Koordinatenebene entsprechen das Element der Fläche (I) $\, .\, .\, .\,\omega\BareRoot{\left(1+\left(\dfrac{\mathrm dz}{\mathrm dx}\right)^2+\left(\dfrac{\mathrm dz}{\mathrm dy}\right)^2\right)}$ und die Menge des magnetischen Fluidums $k\omega\BareRoot{\left(1+\left(\dfrac{\mathrm dz}{\mathrm dx}\right)^2+\left(\dfrac{\mathrm dz}{\mathrm dy}\right)^2\right)}$. Dessen Wirkung auf das Stromelement $i\zeta$ zerfällt also in die drei Kraftkomponenten
\[
\begin{gathered}
ik\zeta\omega.\frac{y}{r^3}\BareRoot{\left(1+\left(\frac{\mathrm dz}{\mathrm dx}\right)^2+\left(\frac{\mathrm dz}{\mathrm dy}\right)^2\right)}\\
-ik\zeta\omega.\frac{x}{r^3}\BareRoot{\left(1+\left(\frac{\mathrm dz}{\mathrm dx}\right)^2+\left(\frac{\mathrm dz}{\mathrm dy}\right)^2\right)}\\
0
\end{gathered}
\]

Nehmen wir jetzt eine zweite Fläche (II), die parallel zu (I) in der unendlich kleinen Entfernung $\varepsilon$ unter ihr liegt. Das heißt, jedem Punkt $x,y,z$ in (I) entspricht in II der Punkt mit den Koordinaten
\[
\begin{gathered}
x+\frac{\varepsilon}{\BareRoot{\left(1+\left(\frac{\mathrm dz}{\mathrm dx}\right)^2+\left(\frac{\mathrm dz}{\mathrm dy}\right)^2\right)}}.\left(\frac{\mathrm dz}{\mathrm dx}\right)\\
y+\frac{\varepsilon}{\BareRoot{\left(1+\left(\frac{\mathrm dz}{\mathrm dx}\right)^2+\left(\frac{\mathrm dz}{\mathrm dy}\right)^2\right)}}.\left(\frac{\mathrm dz}{\mathrm dy}\right)\\
z-\frac{\varepsilon}{\BareRoot{\left(1+\left(\frac{\mathrm dz}{\mathrm dx}\right)^2+\left(\frac{\mathrm dz}{\mathrm dy}\right)^2\right)}}
\end{gathered}
\]

Über diese zweite Fläche sei negatives magnetisches Fluidum so verteilt, dass jeder Flächenteil von II genauso viel negatives Fluidum enthält wie der entsprechende Flächenteil von I positives. Für die Gesamtwirkung der Fluida, die sich auf den einander entsprechenden Elementen von I und II befinden, ergibt sich demnach
\[
\begin{aligned}
\frac{\varepsilon ik\zeta.\omega}{r^5}\left\{3xy.\left(\frac{\mathrm dz}{\mathrm dx}\right)+(3yy-rr)\left(\frac{\mathrm dz}{\mathrm dy}\right)-3yz\right\}&=X\omega\\
-\frac{\varepsilon ik\zeta.\omega}{r^5}\left\{(3xx-rr)\left(\frac{\mathrm dz}{\mathrm dx}\right)+3xy\left(\frac{\mathrm dz}{\mathrm dy}\right)-3xz\right\}&=Y\omega\\
0&=Z\omega
\end{aligned}
\]

Beginnen wir mit der Umformung von $\omega X$, das wir zunächst in folgende Form bringen
\SourceLeaf{632}
\[
\begin{aligned}
\omega X&=\frac{\varepsilon ik\zeta\omega}{r^5}\left\{3xy\left(\frac{\mathrm dz}{\mathrm dx}\right)+(2rr-3xx-3zz)\left(\frac{\mathrm dz}{\mathrm dy}\right)-3yz\right\}\\
&=\frac{\varepsilon ik\zeta\omega}{r^5}\Big\{\left(x\frac{\mathrm ddz}{\mathrm dx\mathrm dy}+\frac{\mathrm dz}{\mathrm dy}\right)rr-3\left(x+z\frac{\mathrm dz}{\mathrm dx}\right)x\frac{\mathrm dz}{\mathrm dy}\\
&\quad-\left(x\frac{\mathrm ddz}{\mathrm dx\mathrm dy}-\frac{\mathrm dz}{\mathrm dy}\right)rr+3\left(y+z\frac{\mathrm dz}{\mathrm dy}\right)\left(x\frac{\mathrm dz}{\mathrm dx}-z\right)\Big\}\\
&=\varepsilon ik\zeta\omega\left\{\frac{\mathrm d\dfrac{x\left(\frac{\mathrm dz}{\mathrm dy}\right)}{r^3}}{\mathrm dx}-\frac{\mathrm d\dfrac{x\left(\frac{\mathrm dz}{\mathrm dx}\right)-z}{r^3}}{\mathrm dy}\right\}
\end{aligned}
\]

Um nun die Gesamtwirkung parallel zur $x$-Achse zu bestimmen, nennen wir die Projektion von I auf die $x,y$-Ebene, also die Gesamtheit aller $\omega$, III. Wir haben daher das Integral $\int\omega X$ über alle $\omega$ zu bestimmen. Schreiben wir $\mathrm dx\mathrm dy$ anstelle von $\omega$, so müssen wir die doppelte Integration von $X\mathrm dx\mathrm dy$ ausführen.

Dabei hat man
\[
\varepsilon ik\zeta\mathrm dy.\left\{\frac{\mathrm d\dfrac{x\left(\frac{\mathrm dz}{\mathrm dy}\right)}{r^3}}{\mathrm dx}.\mathrm dx\right\}
\]
und
\[
\varepsilon ik\zeta\mathrm dx.\left\{\frac{\mathrm d\dfrac{x\left(\frac{\mathrm dz}{\mathrm dx}\right)-z}{r^3}}{\mathrm dy}.\mathrm dy\right\}
\]
getrennt zu betrachten. Für den ersten Ausdruck teilt man III in unendlich viele unendlich schmale Streifen parallel zur $x$-Achse ein, für den zweiten in ähnliche Streifen parallel zur $y$-Achse. Daraus folgt dann sehr leicht, dass das Ganze gleich
\[
\varepsilon ik\zeta\int\frac{x\mathrm dz-z\mathrm dx}{r^3}=\varepsilon ik\zeta\int\frac{x\dfrac{\mathrm dz}{\mathrm ds}-z\dfrac{\mathrm dx}{\mathrm ds}}{r^3}.\mathrm ds
\]
ist, wobei über den gesamten Rand von III integriert wird und man ihn in der Richtung durchläuft, in der III rechts liegt.

Auf ähnliche Weise erhält man für die Summe aller $Y\omega$
\[
\varepsilon ik\zeta\int\frac{y\mathrm dz-z\mathrm dy}{r^3}=\varepsilon ik\zeta\int\frac{y\dfrac{\mathrm dz}{\mathrm ds}-z\dfrac{\mathrm dy}{\mathrm ds}}{r^3}.\mathrm ds
\]
Damit ist in Verbindung mit [Nr. 14] das \textsc{Ampère}sche Gesetz bewiesen.

